\chapter{A Linear-Model Analysis of Placement-Induced Collapse}
\label{app:theory}

\emph{This appendix is a draft analysis to be verified and adapted. It gives a
first-principles explanation, in a linearized model, of why applying the
variance penalty to the context embedding $s_t$ without a co-located
covariance penalty on $s_t$ induces rank collapse---and why that collapse is
invisible to both the training loss and the variance criterion. The argument
reproduces the three empirical signatures reported in
Section~\ref{sec:placement}: a final variance loss of $0$, the largest
$\bar\sigma$, and effective rank $\approx 1$.}

\section{Setup}
Approximate the context encoder and the predictor by linear maps. Let the
input be $x \in \mathbb{R}^{m}$ with $x \sim \mathcal{N}(0, I_m)$, the context
embedding be
\begin{equation}
    s_t = W x, \qquad W \in \mathbb{R}^{d \times m},
\end{equation}
and the predictor be a trainable linear map $s_{\text{pred}} = P\,s_t$ with
$P \in \mathbb{R}^{d \times d}$ (absorbing the action conditioning into $P$ for
this analysis). Write the context-embedding covariance as
\begin{equation}
    \Sigma \;=\; \operatorname{Cov}(s_t) \;=\; W W^{\top} \in \mathbb{R}^{d \times d}.
\end{equation}
The two VICReg-style penalties are, for a chosen tensor $u$ with covariance
$\Sigma_u$,
\begin{align}
    \mathcal{L}_{\text{var}}(u) &= \frac{1}{d}\sum_{j=1}^{d}
        \max\!\big(0,\; \gamma - \sqrt{(\Sigma_u)_{jj}}\big), \\
    \mathcal{L}_{\text{cov}}(u) &= \frac{1}{d}\sum_{i \neq j} (\Sigma_u)_{ij}^{2}.
\end{align}
The variance penalty constrains only the \emph{diagonal} of $\Sigma_u$; the
covariance penalty constrains only the \emph{off-diagonal}.

\section{Original wiring: variance on $s_t$, covariance on $s_{\text{pred}}$}
The objective is
$\mathcal{L} = \mathcal{L}_{\text{distill}}
+ \alpha\,\mathcal{L}_{\text{var}}(s_t) + \beta\,\mathcal{L}_{\text{cov}}(s_{\text{pred}})$.

\paragraph{The covariance term imposes no constraint on the rank of $\Sigma$.}
The covariance of the prediction is $\Sigma_{\text{pred}} = P \Sigma P^{\top}$.
For \emph{any} $\Sigma$ with eigendecomposition $\Sigma = U \Lambda U^{\top}$,
choosing $P = D U^{\top}$ for any diagonal $D$ gives
$\Sigma_{\text{pred}} = D \Lambda D$, which is diagonal, so
$\mathcal{L}_{\text{cov}}(s_{\text{pred}}) = 0$. The predictor can therefore
\emph{absorb} the decorrelation for any context covariance, including a
rank-deficient one. Consequently $\beta\,\mathcal{L}_{\text{cov}}(s_{\text{pred}})$
places no constraint on $\operatorname{rank}(\Sigma)$.

\paragraph{The variance term admits rank-1 solutions.}
The only remaining constraint on $s_t$ is
$\mathcal{L}_{\text{var}}(s_t) = 0 \iff \Sigma_{jj} \ge \gamma^{2}\ \forall j$.
Consider a rank-1 covariance $\Sigma = v v^{\top}$ with
$v \in \mathbb{R}^{d}$ and $|v_j| \ge \gamma$ for all $j$. Then
$\Sigma_{jj} = v_j^{2} \ge \gamma^{2}$, so $\mathcal{L}_{\text{var}}(s_t) = 0$
and the per-dimension standard deviation is $\sqrt{\Sigma_{jj}} = |v_j| \ge \gamma$,
giving $\bar\sigma \ge \gamma$. Its effective rank is exactly $1$.

\paragraph{Conclusion.}
A rank-1 context representation with inflated per-dimension scale is a global
optimum of the original-wiring objective: it drives $\mathcal{L}_{\text{var}}$
and $\mathcal{L}_{\text{cov}}$ to zero and (with a suitable predictor) matches
the distillation target. At this optimum the variance loss is exactly $0$, the
mean standard deviation $\bar\sigma$ is large, and the effective rank is $1$---%
precisely the empirical signature of Table~\ref{tab:placement} (final
variance loss $0.000$, largest $\bar\sigma$, rank $\approx 1$). Because both
the loss and the variance criterion are satisfied, the collapse is invisible
to them.

\section{Shared placement: both penalties on $s_t$}
Now the objective includes $\beta\,\mathcal{L}_{\text{cov}}(s_t)
= \tfrac{\beta}{d}\sum_{i\neq j}\Sigma_{ij}^{2}$, acting on the same covariance
$\Sigma$ that the variance term constrains. A rank-1 $\Sigma = v v^{\top}$ now
incurs an off-diagonal penalty
\begin{equation}
    \sum_{i \neq j} (v_i v_j)^{2}
    = \Big(\sum_j v_j^{2}\Big)^{2} - \sum_j v_j^{4},
\end{equation}
which is large whenever the diagonal constraint $|v_j| \ge \gamma$ holds.
Collapse is therefore no longer optimal. The joint stationary condition drives
the off-diagonals of $\Sigma$ toward $0$ while the diagonals are held at
$\gamma^{2}$, i.e.\ $\Sigma \to \gamma^{2} I_d$, which is full rank $d$. This
is the standard VICReg fixed point, and it explains why co-locating the
covariance term on $s_t$ restores effective rank (Table~\ref{tab:placement},
shared cells).

\section{Why co-location is necessary}
The asymmetry is structural. The variance penalty is a per-coordinate
(diagonal) constraint and is blind to inter-dimensional correlation. Rank
deficiency is an off-diagonal (correlation) property. Preventing it therefore
requires imposing the decorrelation penalty on the \emph{same} representation
whose rank is at stake. Routing that penalty through the trainable predictor
($s_{\text{pred}} = P s_t$) lets $P$ satisfy it for any $\Sigma$, decoupling it
from the rank of $s_t$. Hence a variance penalty without a co-located
covariance penalty inflates the diagonal while leaving the rank unconstrained,
and the collapsed optimum is invisible to the objective---exactly the
phenomenon characterized empirically in this thesis.

\section{Limitations of the analysis}
The argument treats the encoder and predictor as linear and the input as
isotropic Gaussian; it characterizes the optimality of collapsed
\emph{fixed points} rather than proving that gradient descent converges to
them, and it abstracts away batch normalization, the ReLU nonlinearity, and the
EMA target dynamics. A full nonlinear treatment, and a convergence (rather than
optimality) analysis, are left as future work. Nonetheless, the linear model
already predicts the three measured signatures of the collapse and the
effect of co-locating the penalties, which is the intended purpose of this
appendix.
